\documentclass[3p]{elsarticle}

% ── Core packages ─────────────────────────────────────────────────────────────
\usepackage{hyperref}
\usepackage{amssymb}
\usepackage{amsmath}
\usepackage{graphicx}
\graphicspath{{results/}}
\usepackage{booktabs}
\usepackage{siunitx}
\usepackage{makecell}
\usepackage{float}
\usepackage{rotating}
\usepackage{soul}
\usepackage{algorithm}
\usepackage{algorithmic}
\usepackage{subcaption}
\usepackage{threeparttable}
\usepackage{xcolor}
\usepackage{tabularx}
\usepackage[capitalize, noabbrev]{cleveref}
% tikz is not used — flowchart is SVG-based; do not add tikz.

% ── Table caption style ───────────────────────────────────────────────────────
\usepackage[font=normalsize]{caption}
\captionsetup[table]{format=plain, justification=raggedright,
              labelsep=none, singlelinecheck=false, labelfont=bf}

%% Line numbers — uncomment ONLY when submitting to the journal
%% \usepackage{lineno}
%% \linenumbers

% ── Journal name ─────────────────────────────────────────────────────────────
\journal{JOURNAL PLACEHOLDER}

% ─────────────────────────────────────────────────────────────────────────────
\begin{document}

\begin{frontmatter}

% ── Title ─────────────────────────────────────────────────────────────────────
\title{Surgical Duration Prediction with Distilled Clinical Language Models for Operating Room Scheduling}

% ── Authors and affiliations ──────────────────────────────────────────────────
\author[1]{Author One\corref{cor1}}
\author[1]{Author Two}

\affiliation[1]{organization={Department, Faculty, University},
               addressline={Street Address},
               city={City},
               postcode={Postcode},
               state={State/Province},
               country={Country}}

\cortext[cor1]{Corresponding author}
\emailauthor{Author One}{email@university.edu}

% ── Abstract ──────────────────────────────────────────────────────────────────
\begin{abstract}
Accurate prediction of surgical case duration is central to operating-room (OR) efficiency, yet existing machine learning approaches rely almost exclusively on structured Electronic Health Record (EHR) fields and discard the free-text procedure description despite it being the richest pre-operative signal available at the time of booking. Large clinical language models such as Bio-ClinicalBERT can encode this text effectively, but their storage footprints (436\,MB) and inference latencies (21\,ms per case) render them impractical for the resource-constrained hardware common in OR scheduling systems. This paper proposes TinySurgicalBERT, a 0.75\,MB INT8-quantised BERT student model distilled specifically from the surgical procedure lexicon of a 180{,}370-case institutional dataset. A domain-specific Byte-Pair Encoding tokeniser trained on surgical procedure text eliminates out-of-vocabulary tokens, and a combined mean-squared-error and cosine-similarity distillation objective transfers Bio-ClinicalBERT's embedding geometry into a two-layer student with a 128-dimensional hidden representation and a 256-dimensional output embedding. After static INT8 quantisation and ONNX export, the model is evaluated against three encoding baselines — structured-only, SentenceBERT, and Bio-ClinicalBERT — across eight regression models on 180{,}370 surgical cases under five-fold cross-validation with Optuna hyperparameter optimisation. TinySurgicalBERT achieves Mean Absolute Error (MAE) of $26.38 \pm 0.09$ min with XGBoost, statistically equivalent to SentenceBERT ($26.35 \pm 0.07$ min) and Bio-ClinicalBERT ($26.40 \pm 0.11$ min), while requiring $580\times$ and $121\times$ less storage, respectively. All text-augmented encoders reduce MAE by approximately 24\,\% over the structured-only baseline ($34.80 \pm 0.12$ min). End-to-end inference time is $0.484 \pm 0.072$\,ms per case --- $43.7\times$ faster than Bio-ClinicalBERT --- enabling real-time prediction on standard CPU hardware without GPU acceleration. These results demonstrate that domain-specific knowledge distillation can close the gap between predictive accuracy and deployment feasibility in clinical scheduling, making high-quality surgical duration prediction accessible on edge hardware.
\end{abstract}

% ── Graphical abstract (optional — add image path when ready) ─────────────────
\begin{graphicalabstract}
%% \includegraphics[width=\textwidth]{results/graphical_abstract.pdf}
\end{graphicalabstract}

% ── Research highlights (optional — uncomment and fill in before submission) ──
%% \begin{highlights}
%% \item Highlight one
%% \item Highlight two
%% \item Highlight three
%% \item Highlight four
%% \item Highlight five
%% \end{highlights}

% ── Keywords ──────────────────────────────────────────────────────────────────
\begin{keyword}
surgical duration prediction \sep knowledge distillation \sep clinical language models \sep operating room scheduling \sep text encoding \sep BERT quantisation
\end{keyword}

\end{frontmatter}

% ── Abbreviations ─────────────────────────────────────────────────────────────
\section*{Abbreviations}
\begin{description}
    \item[ASA] American Society of Anesthesiologists
    \item[BERT] Bidirectional Encoder Representations from Transformers
    \item[BH] Benjamini--Hochberg
    \item[BPE] Byte-Pair Encoding
    \item[CPU] Central Processing Unit
    \item[CV] Cross-Validation
    \item[EHR] Electronic Health Record
    \item[FDR] False Discovery Rate
    \item[INT8] 8-bit Integer
    \item[MAE] Mean Absolute Error
    \item[MLP] Multi-Layer Perceptron
    \item[NER] Named Entity Recognition
    \item[NLP] Natural Language Processing
    \item[OHE] One-Hot Encoding
    \item[ONNX] Open Neural Network Exchange
    \item[OOV] Out-of-Vocabulary
    \item[OR] Operating Room
    \item[RMSE] Root Mean Square Error
    \item[sMAPE] symmetric Mean Absolute Percentage Error
    \item[TPE] Tree-structured Parzen Estimator
\end{description}

% ── 1. Introduction ───────────────────────────────────────────────────────────
\section{Introduction}\label{sec:introduction}
%% 03_introduction.tex — body only, no preamble

Accurate prediction of surgical case duration is a critical determinant of Operating Room (OR) efficiency, hospital throughput, and patient safety. When scheduled durations deviate substantially from actual case times, the consequences cascade through the entire surgical schedule: underestimates trigger overtime, staff fatigue, and delayed recovery care, while overestimates strand expensive theatre time and force the cancellation or postponement of elective procedures for patients already facing lengthy wait times \cite{oliveira2023systematic,park2025machine}. A systematic review of machine learning approaches to surgical duration prediction found that ML models achieve an average improvement of 25.7\% over traditional scheduling methods, underscoring the magnitude of the gap left by conventional heuristics \cite{park2025machine}.

Operating rooms are among the most expensive resources in any hospital, generating over 40\% of a hospital's revenue while accounting for nearly 30\% of all expenditures; OR time in the United States alone is estimated to cost between \$30 and \$90 billion annually, with each minute of OR time valued at approximately \$46 \cite{al2025comprehensive,park2025machine}. Even modest reductions in scheduling error therefore translate directly into substantial cost savings, improved staff productivity, and greater throughput for waiting patients \cite{wang2022more}.

Existing computational approaches to surgical duration prediction overwhelmingly rely on structured pre-operative Electronic Health Record (EHR) data — surgeon identifier, specialty, patient demographics, and procedure codes — while systematically discarding the free-text procedure description recorded at the time of booking \cite{kwong2025optimizing,park2025machine}. A systematic review of 11 studies published between 2019 and 2024 confirmed that this structured-only paradigm represents the current state of practice, with the best models reducing prediction error by up to 51\% yet still treating the procedure description as non-informative \cite{park2025machine}. Wang and Dexter further demonstrated that even modest reductions in mean absolute predictive error of OR times can yield meaningful increases in labour productivity when coupled with appropriate scheduling adjustments \cite{wang2022more}.

This omission is consequential: the procedure description is the single richest pre-operative signal, encoding the specific surgical plan with a granularity unavailable from any categorical code. Prior work demonstrated that incorporating procedure text into a multimodal prediction model reduces Mean Absolute Error (MAE) by 16\% relative to a structured-only baseline across 180{,}370 retrospective surgical cases, with Sentence-BERT and Bio-ClinicalBERT embeddings achieving the strongest performance \cite{noorchenarboo2026benchmarking}.

Incorporating procedure text effectively requires capable text encoders. Bidirectional Encoder Representations from Transformers (BERT) \cite{devlin2019bert} established state-of-the-art performance across a wide range of Natural Language Processing (NLP) tasks by pre-training deep bidirectional Transformer representations on large unlabelled corpora and fine-tuning with minimal task-specific layers. For sentence-level embedding tasks, Sentence-BERT (SBERT) \cite{reimers2019sentence} extends BERT with siamese and triplet network structures that produce semantically meaningful fixed-size sentence embeddings efficiently — reducing the time to find the most similar pair in a corpus of 10{,}000 sentences from 65 hours with raw BERT to approximately 5 seconds. For the clinical domain, Bio-ClinicalBERT \cite{alsentzer2019publicly} adapts BERT by pre-training on approximately two million clinical notes from the Medical Information Mart for Intensive Care (MIMIC-III) database, improving performance on clinical Named Entity Recognition (NER) and natural language inference tasks. However, both Bio-ClinicalBERT and SBERT carry approximately 110\,M and 22\,M parameters respectively, with disk footprints of 436\,MB and 91\,MB, making them impractical for deployment on the edge or mobile hardware that many hospital scheduling systems rely upon for point-of-care decision support.

Knowledge distillation offers a principled route to compact models that approximate the representational capacity of large teacher networks. Sanh et al.\ demonstrated with DistilBERT \cite{sanh2019distilbert} that a student model trained through a combination of distillation, masked language modelling, and cosine-embedding losses can reduce BERT's parameter count by 40\% and increase inference speed by 60\% while retaining 97\% of language understanding performance. Jiao et al.\ extended this framework with TinyBERT \cite{jiao2020tinybert}, proposing a two-stage distillation strategy — general distillation on a large corpus followed by task-specific fine-tuning — that produces a model 7.5$\times$ smaller and 9.4$\times$ faster than BERT-base on the General Language Understanding Evaluation (GLUE) benchmark. Despite these advances, no prior work has applied knowledge distillation to the surgical procedure text domain, nor evaluated the resulting model within a complete end-to-end surgical duration prediction pipeline that satisfies the latency and storage constraints of real OR scheduling systems.

This paper addresses this gap by proposing TinySurgicalBERT, a lightweight surgical-domain text encoder distilled from Bio-ClinicalBERT and compressed to 0.75\,MB — small enough for deployment on the resource-constrained hardware that OR scheduling systems commonly rely upon. The proposed model is evaluated against Bio-ClinicalBERT and Sentence-BERT across eight regression models on 180{,}370 retrospective surgical cases, demonstrating that domain-specific distillation achieves statistically equivalent predictive accuracy while requiring up to $580\times$ less storage. Together, these results establish that the deployment barrier of large clinical language models can be overcome without sacrificing prediction quality in the surgical scheduling domain. The remainder of this paper is organised as follows. Section~\ref{sec:related_work} reviews prior work on surgical duration prediction and clinical language modelling, identifying the gaps that motivate this work. Section~\ref{sec:methodology} details the data preprocessing, knowledge distillation, and experimental evaluation pipeline. Section~\ref{sec:results} reports results across all encoding and model combinations. Section~\ref{sec:conclusion} draws conclusions and outlines directions for future research.

%%
%% PDFs read: al2025comprehensive.pdf, oliveira2023systematic.pdf, kwong2025optimizing.pdf,
%%            noorchenarboo2026benchmarking.pdf, park2025machine.pdf, wang2022more.pdf,
%%            devlin2019bert.pdf, alsentzer2019publicly.pdf, reimers2019sentence.pdf,
%%            sanh2019distilbert.pdf, jiao2020tinybert.pdf
%% Citations used: oliveira2023systematic, park2025machine, al2025comprehensive, wang2022more,
%%                 kwong2025optimizing, noorchenarboo2026benchmarking, devlin2019bert,
%%                 reimers2019sentence, alsentzer2019publicly, sanh2019distilbert, jiao2020tinybert


% ── 2. Related Work ───────────────────────────────────────────────────────────
\section{Related Work}\label{sec:related_work}
%% 04_related_works.tex — body only, no preamble

Prior work on surgical duration prediction and clinical language modelling spans four thematic areas: machine learning applied to structured Electronic Health Record (EHR) data, artificial intelligence for operating room management and deployment, clinical language models for medical free text, and knowledge distillation methods for producing compact transformer models. Each stream has advanced independently, leaving a critical intersection — domain-specific, deployment-ready text encoding for surgical scheduling — unaddressed.

% ─────────────────────────────────────────────────────────────────────────────
\subsection{Machine Learning for Surgical Duration Prediction}
\label{sec:rw_ml}

Supervised machine learning has been applied to surgical duration prediction across multiple institutional datasets, consistently demonstrating improvements over fixed-schedule heuristics. Bartek et al.\ applied a gradient-boosted ensemble to 46{,}986 retrospective surgical cases and showed that surgeon-specific models outperformed service-level models, with historical case volume per surgeon being the primary determinant of prediction accuracy \cite{bartek2019improving}. Zhao et al.\ evaluated six machine learning algorithms for robot-assisted surgery and found that ensemble methods generalised more reliably across procedure types than linear baselines, particularly for high-variance subspecialties \cite{zhao2019machine}. Martinez et al.\ compared linear regression, support vector machines, regression trees, and bagged trees at a university hospital, confirming that non-linear models outperform linear approaches when procedure-level heterogeneity is high \cite{martinez2021machine}.

Foundational statistical work by Stepaniak et al.\ established that surgical case durations follow a log-normal distribution conditioned on the specific Current Procedural Terminology (CPT) code, anaesthetic type, and individual surgeon, and that modelling these combinations on multi-centre historical data yields more reliable scheduling estimates than institution-wide averages \cite{stepaniak2009modeling,stepaniak2010modeling}. The primary limitation shared across all of these approaches is their exclusive reliance on structured fields: surgeon identifier, specialty, procedure code, and patient demographics constitute the complete feature set, while the free-text procedure description recorded at the time of booking — which encodes the specific surgical plan at a granularity unavailable from any categorical code — is systematically discarded.

% ─────────────────────────────────────────────────────────────────────────────
\subsection{Artificial Intelligence in Operating Room Management}
\label{sec:rw_or}

The broader challenge of operating room management has attracted growing interest from artificial intelligence researchers. Bellini et al.\ surveyed machine learning approaches to OR optimisation and found that decision trees and rule-based models are preferred in clinical deployments due to their interpretability, while gradient-boosted and neural models are largely confined to research settings that lack the governance structures required for clinical adoption \cite{bellini2020artificial}. A subsequent review by the same group identified explainability and transparency as prerequisites for sustained clinical uptake, arguing that end-user trust erodes when models cannot justify their predictions in clinically meaningful terms \cite{bellini2024artificial}. Dios et al.\ demonstrated a decision support system for surgery scheduling deployed across several units of a large Spanish hospital, where the ability to provide an interpretable justification for each scheduling decision was a functional requirement from the outset \cite{dios2015decision}.

While these works advance the case for AI-assisted OR management, they address algorithmic accuracy and interpretability without considering the storage or latency constraints of the hardware on which scheduling systems are actually deployed. None of the reviewed systems incorporates free-text procedure descriptions as a predictive input, and none evaluates whether the proposed models can satisfy the latency requirements of point-of-care scheduling on commodity hardware.

% ─────────────────────────────────────────────────────────────────────────────
\subsection{Clinical Language Models for Medical Text}
\label{sec:rw_clm}

The emergence of transformer-based language models has opened new possibilities for extracting meaning from clinical free text. Li et al.\ developed EhrBERT by fine-tuning BioBERT on 1.5 million unlabelled EHR notes and demonstrated that domain-adapted pre-training substantially improves clinical entity normalisation compared with general-domain BERT, establishing the empirical value of continued pre-training on in-domain clinical corpora \cite{li2019fine}. Turchin et al.\ compared ClinicalBERT, BioBERT, and standard BERT on identifying complex medical concepts in outpatient provider notes, finding that models pre-trained on clinical text consistently outperformed general-corpus models, with clinical domain adaptation accounting for most of the performance gain \cite{turchin2023comparison}. A scoping review by Nunes et al.\ of 47 studies confirmed that domain adaptation through continued pre-training on healthcare text is the dominant paradigm for clinical Natural Language Processing (NLP), with models such as Bio-ClinicalBERT and BioBERT representing the current state of practice across information extraction, relation detection, and clinical coding tasks \cite{nunes2024health}.

Despite their effectiveness, clinical language models carry computational costs that preclude deployment on resource-constrained hardware. The models surveyed by Nunes et al.\ uniformly require 110\,M or more parameters, with on-disk footprints in the hundreds of megabytes, and none has been evaluated under the storage or latency constraints of embedded scheduling systems in perioperative care environments. Critically, no prior work has applied clinical language models to surgical procedure text for the specific purpose of duration prediction.

% ─────────────────────────────────────────────────────────────────────────────
\subsection{Knowledge Distillation and Model Compression}
\label{sec:rw_kd}

Knowledge distillation has become the principled approach to producing compact models that approximate the representational capacity of large teacher networks. Sun et al.\ proposed Patient Knowledge Distillation, which transfers knowledge from intermediate BERT layers rather than only the final output, and showed that a student retaining three to six layers recovers up to 96\,\% of teacher performance on General Language Understanding Evaluation (GLUE) benchmarks while running up to $2\times$ faster at inference \cite{sun2019patient}. Gupta and Agrawal surveyed the complete landscape of text model compression techniques — distillation, pruning, quantisation, and weight sharing — and concluded that knowledge distillation combined with post-training quantisation yields the most favourable accuracy-to-size trade-off across NLP benchmarks \cite{gupta2022compression}. Tang et al.\ extended this analysis to transformer architectures specifically, documenting that layer-reduction distillation strategies consistently outperform width-reduction approaches at equivalent parameter budgets \cite{tang2024survey}.

Bibi et al.\ demonstrated that INT8 static quantisation applied after distillation introduces negligible accuracy degradation on language understanding tasks while halving per-weight storage cost relative to float32 models, confirming the practical viability of the distillation-then-quantise pipeline \cite{bibi2024advances}. Despite these advances in general-domain compression, no prior work has applied knowledge distillation to the surgical procedure text domain, nor has any resulting compressed model been integrated into and evaluated within a complete surgical duration prediction pipeline under the hardware constraints relevant to operating room scheduling.

% ─────────────────────────────────────────────────────────────────────────────
\subsection*{Research Gaps and Distinction of This Work}

Across the four reviewed streams, a consistent and consequential gap emerges: machine learning approaches to surgical duration prediction rely exclusively on structured EHR features and discard the procedure free-text; clinical language models capable of encoding that text are too large for deployment on the resource-constrained hardware common in OR scheduling systems; and knowledge distillation methods that could close this size gap have never been applied to the surgical domain. The consequence is a binary trade-off that no existing method resolves: accurate prediction with a large clinical encoder that cannot be deployed at the point of care, or fast scheduling with a structured-only model that leaves the richest pre-operative signal untapped. This paper resolves the trade-off by distilling Bio-ClinicalBERT on a 180{,}370-case surgical procedure corpus into a 0.75\,MB INT8 student model — TinySurgicalBERT — and evaluating it within a complete five-fold cross-validated prediction pipeline across eight regression models, directly extending the work reviewed in Sections~\ref{sec:rw_ml}--\ref{sec:rw_kd} by occupying the intersection of surgical text encoding, clinical domain adaptation, and deployment-feasible model compression that none of the reviewed works addresses.

\cref{tab:related_summary} summarises the reviewed works and the gap each leaves unaddressed.

\begin{table*}[tp]
\centering
\caption{\\Summary of related works and identified research gaps.}
\label{tab:related_summary}
\begin{threeparttable}
\begin{tabularx}{\linewidth}{@{}lXlX@{}}
\toprule
Reference & Method / Approach & Application Domain & Key Limitation \\
\midrule
\cite{stepaniak2009modeling} & Log-normal CPT-surgeon-anaesthetic model & OR duration prediction & Structured features only; no procedure text \\
\cite{stepaniak2010modeling} & ANOVA surgeon-factor model & OR duration prediction & Requires large per-surgeon history \\
\cite{bartek2019improving}   & Gradient-boosted ensemble, surgeon-specific & OR duration prediction & Structured features only; no procedure text \\
\cite{zhao2019machine}       & Six ML algorithms, robot-assisted surgery & OR duration prediction & Structured features only; no procedure text \\
\cite{martinez2021machine}   & LR / SVM / regression trees comparison & OR duration prediction & Structured features only; no procedure text \\
\cite{dios2015decision}      & DSS for surgery scheduling & OR scheduling deployment & No text encoding; no edge-deployment evaluation \\
\cite{bellini2020artificial} & ML models survey for OR optimisation & OR management & Large models; deployment constraints ignored \\
\cite{bellini2024artificial} & AI explainability review for OR & OR management & Latency and storage not addressed \\
\cite{li2019fine}            & EhrBERT fine-tuned on 1.5\,M EHR notes & Clinical NLP / EHR & 110\,M+ parameters; not deployment-evaluated \\
\cite{turchin2023comparison} & ClinicalBERT / BioBERT / BERT comparison & Clinical NLP & Large models; no surgical scheduling application \\
\cite{nunes2024health}       & Scoping review of 47 healthcare LM studies & Clinical NLP & All models large; edge deployment not evaluated \\
\cite{sun2019patient}        & Patient-KD: intermediate-layer BERT distillation & General NLP & General-domain; no clinical or surgical variant \\
\cite{gupta2022compression}  & Survey: distillation, pruning, quantisation & Text model compression & Not applied to clinical or surgical text \\
\cite{tang2024survey}        & Survey: transformer compression methods & Transformer compression & Not applied to clinical or surgical text \\
\cite{bibi2024advances}      & INT8 quantisation and pruning for NLP & NLP model compression & Not applied to surgical scheduling pipeline \\
\bottomrule
\end{tabularx}
\begin{tablenotes}
\small
\item LR = Linear Regression; SVM = Support Vector Machine; ML = Machine Learning; OR = Operating Room; NLP = Natural Language Processing; KD = Knowledge Distillation; EHR = Electronic Health Record; DSS = Decision Support System.
\end{tablenotes}
\end{threeparttable}
\end{table*}

%% Papers selected: stepaniak2009modeling, stepaniak2010modeling, bartek2019improving,
%%                  zhao2019machine, martinez2021machine, dios2015decision,
%%                  bellini2020artificial, bellini2024artificial, li2019fine,
%%                  turchin2023comparison, nunes2024health, sun2019patient,
%%                  gupta2022compression, tang2024survey, bibi2024advances
%% Papers discarded: liu2021use — Chinese radiology (liver cancer), unrelated to surgical scheduling
%%                   lin2019bert — clinical temporal relation extraction, unrelated task
%%                   acharya2024survey — symbolic KD of LLMs, different paradigm from embedding distillation


% ── 3. Methodology ────────────────────────────────────────────────────────────
\section{Methodology}\label{sec:methodology}
%% 05_methodology.tex — body only, no preamble

Accurate prediction of surgical case duration is a foundational challenge in operating-room scheduling: underestimates cause overtime and patient harm, while overestimates waste costly theatre capacity. Prior approaches rely almost exclusively on structured EHR fields — surgeon identity, specialty, and patient demographics — and discard the free-text procedure description despite it being the single richest predictor available at the time of booking. Exploiting this text is non-trivial: off-the-shelf sentence encoders trained on general or broad clinical corpora exhibit high Out-of-Vocabulary (OOV) rates and poor semantic discrimination for surgical terminology.

This paper addresses this gap with \textbf{TinySurgicalBERT}, a 0.75\,MB INT8-quantised BERT student model distilled specifically from the surgical procedure lexicon of a 180,370-case institutional dataset, designed for real-time surgical-duration prediction on edge and mobile hardware without any cloud dependency. The overall workflow is summarised in Figure~\ref{fig:flowchart}.

% ─────────────────────────────────────────────────────────────────────────────
\subsection{Dataset and Problem Formulation}
\label{sec:dataset}

A retrospective dataset of 194{,}661 surgical cases was obtained from a single tertiary institution over a multi-year window. After applying the data-quality and implausibility filters described in Section~\ref{sec:stage1}, 180{,}370 cases are retained for analysis. Each case is represented by a tuple $(\mathbf{s}, p, y)$ where $\mathbf{s} \in \mathbb{R}^{38}$ is a vector of structured pre-operative features (see Section~\ref{sec:stage1}), $p$ is the free-text procedure description (e.g.\ ``laparoscopic cholecystectomy with intra-operative cholangiogram''), and $y \in \mathbb{R}^{+}$ is the observed surgical duration in minutes. The prediction task is:
\begin{equation}
  \hat{y} = f\!\left(\mathbf{x}\right),
  \qquad
  \mathbf{x} = \bigl[\,\phi(p);\,\mathbf{s}\,\bigr] \in \mathbb{R}^{d},
  \label{eq:task}
\end{equation}
where $\phi(\cdot)$ maps the procedure text to a dense embedding vector and $[\,\cdot\,;\,\cdot\,]$ denotes concatenation. The dimension $d$ depends on the encoding strategy (Section~\ref{sec:stage2}).

% ─────────────────────────────────────────────────────────────────────────────
\subsection{Stage 1 — Data Preprocessing}
\label{sec:stage1}

Raw EHR data contains several systematic artefacts that must be removed before modelling.

\paragraph{Leakage removal.}
Fields that are recorded \emph{after} surgery (e.g.\ actual anaesthesia end time, PACU duration) are excluded from the feature set to prevent target leakage.

\paragraph{Implausibility filter.}
Cases with a recorded duration outside $[0, 2880]$ minutes across all time columns are discarded as implausible scheduling artefacts.

\paragraph{Cyclic time encoding.}
Calendar and clock features (time-of-day, day-of-week, month) are periodic and must not be represented as raw integers. A feature $t$ with period $T$ is mapped to two dimensions via:
\begin{equation}
  t \;\mapsto\; \Bigl(\sin\!\tfrac{2\pi t}{T},\;\cos\!\tfrac{2\pi t}{T}\Bigr),
  \label{eq:cyclic}
\end{equation}
ensuring that, for example, 23:59 and 00:01 are geometrically adjacent.

\paragraph{Categorical One-Hot Encoding (OHE).}
Nominal features (surgeon identifier, specialty, theatre, ASA grade) are one-hot encoded. After encoding and feature selection the structured representation is $\mathbf{s} \in \mathbb{R}^{38}$.

% ─────────────────────────────────────────────────────────────────────────────
\subsection{Stage 2 — Text Encoding and Knowledge Distillation}
\label{sec:stage2}

\subsubsection{Encoding Strategies}

We evaluate four strategies for $\phi(\cdot)$:

\begin{enumerate}
  \item \textbf{Structured Only (baseline).}  $\phi(p) = \mathbf{0}$; the procedure text is discarded. This constitutes a strong structured-only baseline.
  \item \textbf{SentenceBERT.}  A pre-trained general-domain sentence-transformer~\cite{reimers2019sentencebert} maps $p$ to a 384-dimensional embedding. No domain adaptation is performed.
  \item \textbf{Bio-ClinicalBERT.}  A BERT model pre-trained on MIMIC-III clinical notes~\cite{alsentzer2019bioclinicalbert} produces a 768-d pooled embedding via mean-pooling of the last hidden layer.
  \item \textbf{TinySurgicalBERT (ours).}  A lightweight student model (256-d output) obtained via knowledge distillation from Bio-ClinicalBERT, described in Section~\ref{sec:kd}.
\end{enumerate}

\subsubsection{TinySurgicalBERT via Knowledge Distillation}
\label{sec:kd}

\paragraph{Domain-specific tokeniser.}
A Byte-Pair Encoding (BPE) tokeniser is trained on all 180,370 surgical procedure strings in the corpus, with vocabulary size $V = 2500$. This yields an OOV rate of exactly $0\%$ on held-out cases from the same institution, compared with $\sim\!8\%$ for the standard Bio-ClinicalBERT WordPiece vocabulary on the same data.

\paragraph{Architecture.}
The student model uses 2 transformer layers, 4 attention heads, a hidden dimension of 128, an intermediate feed-forward dimension of 256, and a learned linear output projection from 128 to 256 dimensions. The total trainable parameter count is $0.63\text{M}$, compared with $110\text{M}$ for the Bio-ClinicalBERT teacher.

\paragraph{Distillation training.}
Given a procedure string $p$, the teacher embedding $\mathbf{e}_t \in \mathbb{R}^{256}$ is obtained by mean-pooling the final hidden layer of Bio-ClinicalBERT (768-d) and projecting to 256 dimensions via Principal Component Analysis (PCA) fitted on the training corpus. The student embedding $\mathbf{e}_s = \phi_S(p) \in \mathbb{R}^{256}$ is produced by the student's output projection layer (128$\to$256\,d). The student is trained to minimise:
\begin{equation}
  \mathcal{L} \;=\;
    \alpha\,\bigl\lVert \mathbf{e}_s - \mathbf{e}_t \bigr\rVert_2^2
    \;+\;
    (1-\alpha)\,\bigl(1 - \cos(\mathbf{e}_s,\,\mathbf{e}_t)\bigr),
  \label{eq:kd_loss}
\end{equation}
where $\alpha \in (0,1)$ balances mean-squared-error alignment ($\mathcal{L}_{\mathrm{MSE}}$) against cosine similarity maximisation ($\mathcal{L}_{\mathrm{cos}}$). The MSE term encourages magnitude agreement, while the cosine term preserves directional structure regardless of scale.

\paragraph{INT8 quantisation and ONNX export.}
After distillation training, the student is statically quantised to 8-bit Integers (INT8) and exported to ONNX Runtime format, yielding a $0.75\,$MB model file. Inference is performed entirely on CPU using ONNX Runtime, with no GPU or network dependency required at prediction time.

% ─────────────────────────────────────────────────────────────────────────────
\subsection{Feature Assembly}
\label{sec:assembly}

For each encoding strategy the final input vector is constructed by concatenating the text embedding with the structured feature vector:
\begin{equation}
  \mathbf{x} \;=\; \bigl[\,\phi(p);\;\mathbf{s}\,\bigr] \;\in\; \mathbb{R}^{d_\phi + 38},
  \label{eq:concat}
\end{equation}
where $d_\phi \in \{0, 256, 384, 384\}$ for the four encoding strategies respectively. For SentenceBERT and Bio-ClinicalBERT, a fold-wise PCA (fitted on the training fold only) reduces the raw embeddings to 384 principal components; TinySurgicalBERT's 256-dimensional output is used without further reduction.

% ─────────────────────────────────────────────────────────────────────────────
\subsection{Stage 3 — Cross-Validation Protocol}
\label{sec:cv}

To produce unbiased generalisation estimates, all experiments use \textbf{5-fold cross-validation} with random shuffling and a fixed random seed of 42. Each fold comprises 144,296 training cases (80\,\%) and 36,074 validation cases (20\,\%). The same fold assignments are used for all encoding strategies and all downstream models, ensuring that comparisons are paired at the fold level. Statistical tests therefore operate on 5-element performance vectors (one mean per fold), which determines the resolution of significance testing (see Section~\ref{sec:eval}).

% ─────────────────────────────────────────────────────────────────────────────
\subsection{Stage 4 — Hyperparameter Optimisation and Downstream Models}
\label{sec:hpo}

\paragraph{Model family.}
Eight regression models are evaluated: Linear Regression, Ridge, Lasso, ElasticNet, Random Forest, XGBoost, LightGBM, and a Multi-Layer Perceptron (MLP). This selection spans the spectrum from interpretable linear baselines to high-capacity non-linear ensembles.

\paragraph{Optuna TPE search.}
Each model is independently tuned using \textbf{Optuna}~\cite{akiba2019optuna} with the Tree-structured Parzen Estimator (TPE) sampler. A budget of \textbf{20 trials} is used per (encoding~$\times$~model) combination, with the first 5 trials drawn randomly to initialise the surrogate. Hyperparameter spaces are defined per model: regularisation strength for linear models; tree depth, number of estimators, and learning rate for ensemble methods; hidden size, number of layers, and dropout rate for the MLP. The objective metric for TPE is Mean Absolute Error (MAE) on the inner cross-validation fold.

% ─────────────────────────────────────────────────────────────────────────────
\subsection{Evaluation Metrics}
\label{sec:eval}

All metrics are computed per fold and reported as mean~$\pm$ standard deviation across the five folds. The four primary metrics are:
\begin{align}
  \mathrm{MAE}   &= \frac{1}{n}\sum_{i=1}^{n}\lvert y_i - \hat{y}_i\rvert,
                 \label{eq:mae}\\[4pt]
  \mathrm{RMSE}  &= \sqrt{\frac{1}{n}\sum_{i=1}^{n}(y_i - \hat{y}_i)^2},
                 \label{eq:rmse}\\[4pt]
  \mathrm{sMAPE} &= \frac{200}{n}\sum_{i=1}^{n}
                    \frac{\lvert y_i - \hat{y}_i\rvert}{\lvert y_i\rvert+\lvert\hat{y}_i\rvert},
                 \label{eq:smape}\\[4pt]
  R^{2}          &= 1 - \frac{\sum_{i}(y_i-\hat{y}_i)^2}{\sum_{i}(y_i-\bar{y})^2}.
                 \label{eq:r2}
\end{align}
MAE and Root Mean Square Error (RMSE) are in units of minutes; the symmetric Mean Absolute Percentage Error (sMAPE) is scale-invariant and symmetric, avoiding the blow-up of standard MAPE near zero; $R^{2}$ measures the fraction of duration variance explained by the model.

\paragraph{Statistical comparison.}
Pairwise model comparisons use the \textbf{Wilcoxon signed-rank test} on the 5 per-fold metric vectors (two-sided, $\alpha = 0.05$). Because $n = 5$ folds is the minimum sample size at which the test can achieve a $p$-value below $0.0625$, all $p$-values from this test are reported alongside the raw fold-level differences; multiple comparisons across encoding strategies are corrected using the \textbf{Benjamini–Hochberg} False Discovery Rate (FDR-BH) procedure.


%% Flowchart is referenced inside 03_methodology as Figure~\ref{fig:flowchart}
\begin{figure*}[tp]
\centering
\input{flowchart/06_flowchart}
\caption{Overall workflow of the proposed methodology.}
\label{fig:flowchart}
\end{figure*}

% ── 4. Results and Discussion ─────────────────────────────────────────────────
\section{Results and Discussion}\label{sec:results}
%% 07_results.tex — body only, no preamble

This section reports results across three dimensions: predictive accuracy across all 32 encoding–model combinations, deployment efficiency of the three text encoders, and pairwise statistical comparisons. The central findings are: (i) all three text-augmented encoding strategies reduce MAE by approximately 24\,\% over the structured-only baseline; (ii) TinySurgicalBERT achieves accuracy statistically equivalent to Bio-ClinicalBERT and SentenceBERT while requiring up to $580\times$ less on-disk storage; and (iii) TinySurgicalBERT runs at $0.484$\,ms per case — $43.7\times$ faster than Bio-ClinicalBERT — satisfying real-time OR scheduling requirements on standard CPU hardware without GPU acceleration. All numerical results are derived directly from the experimental database without post-hoc modification.

% ─────────────────────────────────────────────────────────────────────────────
\subsection{Dataset Description}
\label{sec:dataset_desc}

A retrospective dataset of 194{,}661 surgical cases was obtained from a single tertiary institution. After applying the data-quality and implausibility filters described in Section~\ref{sec:stage1}, 180{,}370 cases are retained for modelling. The prediction target is the observed procedure duration in minutes, with a mean of $95.0$ minutes, a standard deviation of $92.1$ minutes, and a median of $66.0$ minutes, reflecting a strongly right-skewed distribution characteristic of surgical workload data. Patient age is available for $153{,}669$ cases (85.2\,\% completion; mean $56.6 \pm 17.6$ years), and body mass index for $119{,}193$ cases (66.1\,\% completion; mean $27.8 \pm 9.1\,\text{kg}\,\text{m}^{-2}$). Cases span 13 surgical services, with Orthopaedic Surgery (27.0\,\%), General Surgery (14.7\,\%), and Obstetrics/Gynaecology (13.6\,\%) comprising the three largest groups. Anaesthetic type is predominantly general (84.6\,\%), with neuraxial (7.8\,\%), sedation (3.2\,\%), and regional (2.7\,\%) procedures also represented. Descriptive statistics for the principal pre-operative features are summarised in \cref{tab:dataset_desc}.

\begin{table}[tp]
\centering
\caption{\\Descriptive statistics for the primary pre-operative features. Continuous features report minimum, maximum, mean, and standard deviation computed from all non-missing records. Categorical and binary features report the number of distinct levels. BMI and age contain missing values as indicated.}
\label{tab:dataset_desc}
\begin{threeparttable}
\begin{tabular}{llrrrr}
\toprule
Feature & Type & Min & Max & Mean & Std \\
\midrule
Procedure duration (min)\tnote{a} & Continuous & 0.0 & 1{,}535.0 & 95.0 & 92.1 \\
Scheduled duration (min) & Continuous & 15.0 & 1{,}870.0 & 151.3 & 105.8 \\
Patient age (years)\tnote{b} & Continuous & 18.0 & 103.1 & 56.6 & 17.6 \\
Body mass index (kg m$^{-2}$)\tnote{c} & Continuous & 6.1 & 199.5 & 27.8 & 9.1 \\
OR team size & Continuous & 1.0 & 44.0 & 10.0 & 3.1 \\
Scheduled start hour & Continuous & 0.0 & 18.0 & 10.9 & 2.5 \\
ASA physical status score & Ordinal & 1 & 5 & — & — \\
Patient sex & Binary & \multicolumn{4}{l}{54.0\,\% male ($n = 96{,}994$)} \\
Case service & Categorical & \multicolumn{4}{l}{13 levels} \\
Anaesthetic type & Categorical & \multicolumn{4}{l}{5 levels} \\
Procedure description & Text & \multicolumn{4}{l}{free-text; encoded by $\phi(\cdot)$} \\
\bottomrule
\end{tabular}
\begin{tablenotes}
\small
\item[a] Prediction target; right-skewed (median $66.0$ min).
\item[b] Available for $153{,}669$ cases (85.2\,\%).
\item[c] Available for $119{,}193$ cases (66.1\,\%).
\end{tablenotes}
\end{threeparttable}
\end{table}

% ─────────────────────────────────────────────────────────────────────────────
\subsection{Experimental Setup}
\label{sec:exp_setup}

All experiments use 5-fold cross-validation with a fixed random seed of 42, as described in Section~\ref{sec:cv}, producing training and validation sets of $144{,}296$ and $36{,}074$ cases per fold, respectively. For each of the 32 encoding-model combinations (4 encoding strategies $\times$ 8 regression models), Optuna with the Tree-structured Parzen Estimator (TPE) sampler allocates 20 trials per combination, with the first 5 trials drawn uniformly at random to initialise the surrogate model. Hyperparameter search is conducted independently per fold to avoid information leakage from validation data into the tuning process. The objective metric throughout is Mean Absolute Error (MAE) on the held-out fold. Optimal hyperparameters obtained for the TinySurgicalBERT encoding from a representative fold are reported in \cref{tab:hyperparams}; values are consistent across folds (coefficient of variation $< 5\%$ for continuous parameters).

\begin{table}[tp]
\centering
\caption{\\Optimal hyperparameters identified by Optuna TPE search for the TinySurgicalBERT encoding (fold 0). Search space bounds and a brief description of each parameter are provided. Parameters for linear models (Ridge, Lasso, ElasticNet) are omitted for brevity; regularisation strengths converge to $\alpha \approx 10^{-3}$--$10^{-1}$.}
\label{tab:hyperparams}
\begin{threeparttable}
\resizebox{\linewidth}{!}{%
\begin{tabular}{llrll}
\toprule
Model & Parameter & Optimal value & Search range & Description \\
\midrule
XGBoost & \texttt{n\_estimators}   & 440   & [50, 500]          & Number of boosting rounds \\
        & \texttt{learning\_rate}  & 0.056 & [$10^{-3}$, 0.3]   & Step size shrinkage \\
        & \texttt{max\_depth}      & 8     & [3, 10]            & Maximum tree depth \\
        & \texttt{subsample}       & 0.90  & [0.5, 1.0]         & Row sub-sampling ratio \\
\midrule
LightGBM & \texttt{n\_estimators}  & 482   & [50, 500]          & Number of boosting rounds \\
         & \texttt{learning\_rate} & 0.078 & [$10^{-3}$, 0.3]   & Step size shrinkage \\
         & \texttt{max\_depth}     & 8     & [3, 12]            & Maximum tree depth \\
         & \texttt{num\_leaves}    & 124   & [20, 150]          & Maximum number of leaves \\
         & \texttt{subsample}      & 0.69  & [0.5, 1.0]         & Row sub-sampling ratio \\
\midrule
MLP      & \texttt{n\_layers}      & 2     & [1, 4]             & Number of hidden layers \\
         & \texttt{activation}     & elu   & \{relu, elu, tanh\} & Hidden activation function \\
         & \texttt{lr}             & 0.0085& [$10^{-4}$, $10^{-2}$] & AdamW learning rate \\
         & \texttt{units\_1}       & 141   & [64, 256]          & Neurons in first hidden layer \\
\midrule
Random Forest & \texttt{n\_estimators}  & 126 & [50, 300]        & Number of trees \\
              & \texttt{max\_depth}     & 10  & [5, 20]          & Maximum tree depth \\
              & \texttt{max\_features}  & 0.50 & [0.2, 1.0]      & Features considered per split \\
\bottomrule
\end{tabular}}
\begin{tablenotes}
\small
\item All values from fold 0; variation across folds is negligible (coefficient of variation $< 5\%$ for all continuous parameters).
\end{tablenotes}
\end{threeparttable}
\end{table}

% ─────────────────────────────────────────────────────────────────────────────
\subsection{Predictive Performance}
\label{sec:perf}

\cref{fig:pred_metrics} reports mean and standard deviation of the four primary metrics across all 32 encoding-model combinations. The clearest finding is the division between tree-based ensembles and linear models: XGBoost and LightGBM consistently achieve the lowest MAE and sMAPE regardless of encoding strategy, while linear models (Ridge, Lasso, ElasticNet, Linear Regression) plateau near MAE $\approx 35$--47 min depending on encoding. This pattern indicates that surgical duration is governed by non-linear interactions among the pre-operative features that gradient-boosted trees exploit effectively but that linear predictors cannot capture.

The most practically important finding concerns the relative standing of the four encoding strategies for non-linear models. The structured-only baseline, which discards the procedure description entirely, produces MAE of $34.80 \pm 0.12$ min for XGBoost and $34.79 \pm 0.17$ min for LightGBM. Adding any form of text embedding reduces MAE to the $26$--$27$ minute range across all three text-augmented strategies, representing an absolute reduction of approximately $8.4$ minutes (24\,\% improvement) for the best model. This reduction is consistent across all non-linear models and all folds, indicating that the effect is robust rather than artefactual.

Among the three text-encoding strategies, TinySurgicalBERT achieves $\text{MAE} = 26.38 \pm 0.09$ min and $R^{2} = 0.8543 \pm 0.0051$ with XGBoost and $\text{MAE} = 26.43 \pm 0.10$ min with LightGBM, statistically indistinguishable from SentenceBERT ($\text{MAE} = 26.35 \pm 0.07$, $R^{2} = 0.8542 \pm 0.0049$) and Bio-ClinicalBERT ($\text{MAE} = 26.40 \pm 0.11$, $R^{2} = 0.8539 \pm 0.0054$). This parity is achieved with a $580\times$ reduction in model size relative to Bio-ClinicalBERT (0.75\,MB vs.\ 436\,MB) and a $121\times$ reduction relative to SentenceBERT (0.75\,MB vs.\ 91\,MB). The MLP downstream model shows notably higher fold-to-fold variability under TinySurgicalBERT ($\text{MAE} = 26.95 \pm 0.26$ min) compared to tree ensembles ($\pm 0.09$--$0.10$ min), consistent with the sensitivity of neural networks to random weight initialisation.

\begin{figure}[tp]
  \centering
  \includegraphics[width=\linewidth]{pred_legend}
  \par\medskip
  \begin{subfigure}[b]{0.48\textwidth}
    \includegraphics[width=\textwidth]{pred_mae}
    \caption{MAE (min, $\downarrow$)}
    \label{fig:pred_mae}
  \end{subfigure}\hfill
  \begin{subfigure}[b]{0.48\textwidth}
    \includegraphics[width=\textwidth]{pred_mse}
    \caption{RMSE (min, $\downarrow$)}
    \label{fig:pred_rmse}
  \end{subfigure}
  \par\medskip
  \begin{subfigure}[b]{0.48\textwidth}
    \includegraphics[width=\textwidth]{pred_smape}
    \caption{sMAPE (\%, $\downarrow$)}
    \label{fig:pred_smape}
  \end{subfigure}\hfill
  \begin{subfigure}[b]{0.48\textwidth}
    \includegraphics[width=\textwidth]{pred_r2}
    \caption{$R^{2}$ ($\uparrow$)}
    \label{fig:pred_r2}
  \end{subfigure}
  \caption{Predictive performance across all encoding strategies and downstream models (mean $\pm$ standard deviation over five cross-validation folds). The shared legend above identifies each encoding. Lower is better for MAE, RMSE, and sMAPE; higher is better for $R^{2}$.}
  \label{fig:pred_metrics}
\end{figure}

% ─────────────────────────────────────────────────────────────────────────────
\subsection{Model Size and Inference Efficiency}
\label{sec:runtime}

Having established that all three text-augmented encoders deliver equivalent predictive accuracy, the question becomes whether TinySurgicalBERT can also meet the storage and latency constraints of real OR scheduling systems. TinySurgicalBERT is deployed as an INT8 ONNX model; SentenceBERT and Bio-ClinicalBERT are loaded as fp32 HuggingFace safetensors models. \cref{fig:efficiency} compares the three text encoders on two deployment-critical dimensions: on-disk model file size and end-to-end per-case inference time (encoder forward pass plus downstream XGBoost prediction).

TinySurgicalBERT (INT8 ONNX, $0.75$\,MB) is $580\times$ smaller than Bio-ClinicalBERT ($436$\,MB) and $121\times$ smaller than SentenceBERT ($91$\,MB), as shown in \cref{fig:model_size}. The compression is achieved through the combination of a compact two-layer student architecture (0.63\,M parameters vs.\ 110\,M in Bio-ClinicalBERT) and static INT8 quantisation, which halves the per-weight storage cost relative to a float32 model of the same depth.

End-to-end inference times, measured under single-item CPU inference to reflect point-of-care scheduling conditions, are presented in \cref{fig:infer_time} and \cref{tab:efficiency}. Across $N = 30$ repeated passes over 200 procedure descriptions, TinySurgicalBERT requires $0.484 \pm 0.072$\,ms per case, compared with $2.570 \pm 0.136$\,ms for SentenceBERT ($5.3\times$ slower) and $21.190 \pm 0.111$\,ms for Bio-ClinicalBERT ($43.7\times$ slower). Both differences are statistically significant ($p < 0.001$, two-sided paired Wilcoxon signed-rank, FDR-BH corrected). Downstream XGBoost prediction contributes less than $0.005$\,ms per case for all three encoders and is negligible relative to the encoding step. These results confirm that TinySurgicalBERT meets real-time scheduling requirements on commodity CPU hardware, a deployment scenario that neither SentenceBERT nor Bio-ClinicalBERT can satisfy without GPU acceleration or batch pre-computation.

\begin{figure}[tp]
  \centering
  \includegraphics[width=\linewidth]{efficiency_legend}
  \par\medskip
  \begin{subfigure}[b]{0.48\textwidth}
    \includegraphics[width=\textwidth]{model_size}
    \caption{Encoder model file size (MB, log scale)}
    \label{fig:model_size}
  \end{subfigure}\hfill
  \begin{subfigure}[b]{0.48\textwidth}
    \includegraphics[width=\textwidth]{infer_time}
    \caption{End-to-end inference time per case (ms, log scale)}
    \label{fig:infer_time}
  \end{subfigure}
  \caption{Deployment efficiency of the three text encoders. Model size is the on-disk file size of the deployed encoder artefact (INT8 ONNX for TinySurgicalBERT; fp32 HuggingFace safetensors for the others). Inference time is the total per-case time from raw procedure text to predicted duration, measured under single-item CPU inference (encoding $+$ XGBoost prediction; downstream contribution $<$0.005\,ms). Structured Only is excluded as its predictive inferiority was established in Section~\ref{sec:perf}.}
  \label{fig:efficiency}
\end{figure}

%% results/efficiency_table.tex — full-width deployment comparison table
%% Deployment format column removed (mentioned in prose); resizebox added for fit.

\begin{table*}[tp]
\centering
\caption{\\Deployment profile of the three text encoders. Architecture columns
  describe the deployed model. Inference time is mean\,$\pm$\,SD over $N=30$
  repeated passes on 200 surgical procedure descriptions under single-item CPU
  inference. Compression and speed-up are relative to TinySurgicalBERT (ours).
  {\color{green!60!black}$\checkmark$}\,=\,TinySurgicalBERT significantly faster
  (two-sided paired Wilcoxon signed-rank, Benjamini--Hochberg FDR corrected).}
\label{tab:efficiency}
\begin{threeparttable}
{\renewcommand{\tabcolsep}{4pt}%
\resizebox{\textwidth}{!}{%
\begin{tabular}{lccccrrrr}
\toprule
Encoder &
  \makecell{Trans-\\former\\layers} &
  \makecell{Atten-\\tion\\heads} &
  \makecell{Hidden\\dim} &
  \makecell{Params\\(M)} &
  \makecell{Size\\(MB)} &
  \makecell{Compres-\\sion\\($\times$)} &
  \makecell{Inference\\time\\(ms/case)} &
  \makecell{Speed-\\up\\($\times$)} \\
\midrule
TinySurgicalBERT (ours) & 2  & 4  & 128 & 0.63  & $0.752$  & ---  & $0.484 \pm 0.072$ & --- \\
SentenceBERT            & 6  & 12 & 384 & 22.7  & $90.87$  & $121$  & $2.570 \pm 0.136$\;{\color{green!60!black}$\checkmark^{***}$} & $5.3$ \\
Bio-ClinicalBERT        & 12 & 12 & 768 & 110.0 & $435.76$ & $580$  & $21.190 \pm 0.111$\;{\color{green!60!black}$\checkmark^{***}$} & $43.7$ \\
\bottomrule
\end{tabular}}
}% end renewcommand group
\begin{tablenotes}
\small
\item $^{***}p < 0.001$ (two-sided paired Wilcoxon signed-rank,
  Benjamini--Hochberg FDR corrected, $N=30$).
\item Downstream XGBoost prediction contributes ${<}\,0.005$\,ms/case for all
  encoders (negligible).
\end{tablenotes}
\end{threeparttable}
\end{table*}


% ─────────────────────────────────────────────────────────────────────────────
\subsection{Statistical Comparison}
\label{sec:stat}

Pairwise comparisons between TinySurgicalBERT and each baseline encoding are evaluated using the two-sided Wilcoxon signed-rank test on the five per-fold metric vectors, with multiple comparisons corrected by the Benjamini--Hochberg False Discovery Rate (FDR-BH) procedure across all 12 test pairs (3 baselines $\times$ 4 metrics). Two distinct patterns emerge. Against the structured-only baseline, all five per-fold MAE differences favour TinySurgicalBERT (mean difference $-8.42$ min, $W = 0$, $p = 0.0625$ for XGBoost), with the same unanimity holding for RMSE, sMAPE, and $R^{2}$: $W = 0$ in every case, the strongest evidence achievable with $n = 5$ folds. Against SentenceBERT and Bio-ClinicalBERT, per-fold differences are small and unsystematic ($W \geq 4$, $p \geq 0.44$), providing no evidence of any performance gap. This pattern of maximal evidence of superiority over the structured-only baseline combined with no evidence of difference against the large encoders is the key result: TinySurgicalBERT performs at the level of models $121\times$--$580\times$ its size. \cref{tab:stat} reports the full pairwise comparison.

\begin{table}[tp]
\centering
\caption{\\Pairwise statistical comparison of TinySurgicalBERT + XGBoost (reference) against baseline encoding strategies + XGBoost. Values show mean $\pm$ standard deviation across five cross-validation folds. Significance markers indicate the outcome of the two-sided Wilcoxon signed-rank test with Benjamini--Hochberg FDR correction ($n = 5$ folds): {\color{green!55!black}$\checkmark$} = reference significantly better; {\color{red!70!black}$\times$} = reference significantly worse; blank = no significant difference. Column arrows indicate the direction of improvement.}
\label{tab:stat}
\begin{threeparttable}
\resizebox{\linewidth}{!}{%
\begin{tabular}{lcccc}
\toprule
Method &
  MAE $\downarrow$ (min) &
  RMSE $\downarrow$ (min) &
  sMAPE $\downarrow$ (\%) &
  $R^{2}$ $\uparrow$ \\
\midrule
TinySurgicalBERT + XGBoost (ours) &
  $26.38 \pm 0.09$ &
  $41.87 \pm 0.77$ &
  $21.61 \pm 0.15$ &
  $0.8543 \pm 0.0051$ \\
\midrule
Structured Only + XGBoost &
  $34.80 \pm 0.12$ &
  $51.77 \pm 0.67$ &
  $28.84 \pm 0.12$ &
  $0.7773 \pm 0.0050$ \\
SentenceBERT + XGBoost &
  $26.35 \pm 0.07$ &
  $41.89 \pm 0.75$ &
  $21.56 \pm 0.06$ &
  $0.8542 \pm 0.0049$ \\
Bio-ClinicalBERT + XGBoost &
  $26.40 \pm 0.11$ &
  $41.93 \pm 0.81$ &
  $21.62 \pm 0.09$ &
  $0.8539 \pm 0.0054$ \\
\bottomrule
\end{tabular}}
\begin{tablenotes}
\small
\item $^{*}p < 0.05$;\; $^{**}p < 0.01$;\; $^{***}p < 0.001$ (two-sided Wilcoxon signed-rank, Benjamini--Hochberg FDR corrected).
\item With $n = 5$ folds the minimum achievable two-sided $p$-value is $0.0625$; no comparison therefore reaches the $\alpha = 0.05$ threshold. Against Structured Only, all fold-level differences are consistently signed ($W = 0$, $p = 0.0625$), the strongest evidence achievable at this sample size.
\item $\downarrow$ lower is better;\; $\uparrow$ higher is better.
\end{tablenotes}
\end{threeparttable}
\end{table}


% ── 5. Conclusion ─────────────────────────────────────────────────────────────
\section{Conclusion}\label{sec:conclusion}
This work addressed the challenge of deploying clinical language models for surgical duration prediction on the resource-constrained hardware common in operating-room scheduling systems. Existing approaches either rely on structured Electronic Health Record (EHR) features alone, discarding the information-rich procedure free-text, or require large clinical encoders such as Bio-ClinicalBERT that are impractical for edge deployment. TinySurgicalBERT was proposed as a principled resolution: a domain-specific BERT student model distilled from Bio-ClinicalBERT on a 180{,}370-case surgical procedure corpus, quantised to INT8, and exported to ONNX Runtime format, yielding a 0.75\,MB encoder that operates entirely on CPU without any GPU or network dependency at inference time.

Evaluated across 32 encoding--model combinations on 180{,}370 surgical cases under five-fold cross-validation, TinySurgicalBERT achieved MAE of $26.38 \pm 0.09$ min with XGBoost --- statistically equivalent to SentenceBERT ($26.35 \pm 0.07$ min) and Bio-ClinicalBERT ($26.40 \pm 0.11$ min) --- while requiring $580\times$ and $121\times$ less on-disk storage, respectively. All three text-augmented encoders improved upon the structured-only baseline by approximately 24\,\% in MAE (from $34.80 \pm 0.12$ to $26.38 \pm 0.09$ min with XGBoost), establishing that procedure free-text provides independent predictive signal beyond structured EHR fields alone. TinySurgicalBERT further demonstrated inference at $0.484 \pm 0.072$\,ms per case --- $43.7\times$ faster than Bio-ClinicalBERT --- confirming suitability for real-time point-of-care scheduling on standard CPU hardware.

Several limitations of this study should be acknowledged. The dataset originates from a single tertiary institution, and both the domain-specific tokeniser and the distillation objective were tuned on its particular surgical vocabulary; performance may differ at sites with distinct procedural nomenclature or scheduling conventions, and institution-specific re-distillation may be required. The five-fold cross-validation protocol limits the statistical resolution of pairwise comparisons: with $n = 5$ folds, the minimum achievable two-sided Wilcoxon $p$-value is $0.0625$, which prevents formal equivalence testing at $\alpha = 0.05$ even when observed differences are negligible in magnitude. Additionally, the evaluation is restricted to pre-operative features available at the time of booking; intra-operative signals that become available in real time were deliberately excluded to preserve the pre-operative prediction setting but may offer further predictive gains in adaptive scheduling frameworks.

Three directions for future work emerge from these findings. First, external validation across multiple institutions with differing procedural vocabularies would establish the generalisability of TinySurgicalBERT and clarify when institution-specific distillation is necessary. Second, the confirmed equivalence between TinySurgicalBERT and the large encoders motivates investigation of more aggressive compression strategies --- structured pruning, low-rank decomposition, or 4-bit quantisation --- to assess whether further size reductions below 0.75\,MB are achievable without accuracy degradation. Third, prospective deployment within a live OR scheduling system, with downstream measurement of utilisation, overtime rates, and case cancellations, would provide direct evidence of operational benefit and close the gap between predictive accuracy on retrospective data and real-world clinical impact.

% ── CRediT authorship contribution statement ──────────────────────────────────
\section*{CRediT authorship contribution statement}
\textbf{Author One}: Conceptualization, Methodology, Data analysis and curation,
Software, Writing -- original draft and editing.
\textbf{Author Two}: Conceptualization, Methodology, Supervision,
Writing -- review and editing, Guidance, and Oversight throughout the research process.

% ── Declaration of competing interest ─────────────────────────────────────────
\section*{Declaration of competing interest}
The authors declare that they have no known competing financial interests or
personal relationships that could have appeared to influence the work reported
in this paper.

% ── Acknowledgements ──────────────────────────────────────────────────────────
\section*{Acknowledgement}
%% Fill in funding sources and support before submission.

% ── Data availability ─────────────────────────────────────────────────────────
\section*{Data availability}
%% Choose one option and remove the others before submission:
%% Option A — data is open:
%%   The data used in this study are publicly available at [repository/DOI].
%% Option B — data is proprietary:
%%   The data used in this study is proprietary and subject to confidentiality
%%   agreements; therefore, it cannot be shared.
%% Option C — data available on request:
%%   The data that support the findings of this study are available from the
%%   corresponding author upon reasonable request.

% ── Bibliography ──────────────────────────────────────────────────────────────
\bibliographystyle{unsrt}
\bibliography{overleaf/references}

\end{document}
